\documentclass[12pt]{article}
\usepackage[spanish]{babel}
\usepackage{amsmath}
\usepackage{amssymb}
\usepackage{booktabs}
\usepackage{geometry}
\geometry{a4paper, margin=1.5cm}

\title{Bondad de Ajuste en Modelos de Regresión Lineal}
\author{Machine Learning - Evaluación de Modelos}
\date{}

\begin{document}

\maketitle



En el contexto de modelos de regresión lineal, la \textbf{bondad de ajuste} se refiere a la capacidad del modelo para capturar adecuadamente la relación entre las variables independientes y la variable dependiente. Un buen ajuste indica que el modelo explica una proporción significativa de la variabilidad en los datos.



\section{Modelo de Regresión Lineal}

El modelo de regresión lineal múltiple se expresa como:

\[
Y = X\beta + \varepsilon
\]

donde:
\begin{itemize}
\item $Y$ es el vector de observaciones ($n \times 1$)
\item $X$ es la matriz de diseño ($n \times p$)
\item $\beta$ es el vector de parámetros ($p \times 1$)
\item $\varepsilon$ es el vector de errores ($n \times 1$)
\end{itemize}
%%%%%%%%%%%%%%%%%%%%%%%%%%%%%%%%%%%%%%%%%%%%%%%%%%%%%%%%%%%%%%%
\section{Descomposición de la Variabilidad}



La descomposición de la variabilidad es fundamental para entender cómo el modelo de regresión explica la dispersión de los datos. Esta descomposición se basa en el análisis de varianza (ANOVA) y proporciona la base teórica para el cálculo del coeficiente de determinación $R^2$.


\subsection{Suma Total de Cuadrados (SST)}

Representa la variabilidad total presente en la variable dependiente $Y$:

\[
\text{SST} = \sum_{i=1}^{n}(y_i - \bar{y})^2 = \mathbf{Y}^T\mathbf{Y} - n\bar{y}^2
\]

donde $\bar{y} = \frac{1}{n}\sum_{i=1}^{n}y_i$ es la media muestral de $Y$.

\subsection{Suma de Cuadrados de la Regresión (SSR)}

Mide la variabilidad explicada por el modelo:

\[
\text{SSR} = \sum_{i=1}^{n}(\hat{y}_i - \bar{y})^2 = \hat{\boldsymbol{\beta}}^T\mathbf{X}^T\mathbf{Y} - n\bar{y}^2
\]

\subsection{Suma de Cuadrados de los Errores (SSE)}

Cuantifica la variabilidad no explicada por el modelo:

\[
\text{SSE} = \sum_{i=1}^{n}(y_i - \hat{y}_i)^2 = \mathbf{Y}^T\mathbf{Y} - \hat{\boldsymbol{\beta}}^T\mathbf{X}^T\mathbf{Y}
\]

\textbf{Demostración de la Descomposición}

\begin{align*}
\text{SST} &= \sum_{i=1}^{n}(y_i - \bar{y})^2 \\
&= \sum_{i=1}^{n}[(y_i - \hat{y}_i) + (\hat{y}_i - \bar{y})]^2 \\
&= \sum_{i=1}^{n}(y_i - \hat{y}_i)^2 + \sum_{i=1}^{n}(\hat{y}_i - \bar{y})^2 + 2\sum_{i=1}^{n}(y_i - \hat{y}_i)(\hat{y}_i - \bar{y})
\end{align*}

El término cruzado se anula:

\begin{align*}
\sum_{i=1}^{n}(y_i - \hat{y}_i)(\hat{y}_i - \bar{y}) &= \sum_{i=1}^{n}\varepsilon_i(\hat{y}_i - \bar{y}) \\
&= \boldsymbol{\varepsilon}^T(\mathbf{\hat{Y}} - \bar{y}\mathbf{1}) \\
&= 0 \quad \text{(por ortogonalidad de los residuales)}
\end{align*}

Por lo tanto:

\[
\text{SST} = \text{SSE} + \text{SSR}
\]

\textbf{Grados de Libertad}

\begin{table}[h]
\centering
\caption{Descomposición de grados de libertad}
\begin{tabular}{lcccc}
\toprule
Fuente de Variación & Suma de Cuadrados & Grados de Libertad & Cuadrado Medio & Esperanza del Cuadrado Medio \\
\midrule
Regresión & $\text{SSR}$ & $p-1$ & $\text{MSR} = \frac{\text{SSR}}{p-1}$ & $\sigma^2 + \frac{\boldsymbol{\beta}^T\mathbf{X}^T\mathbf{X}\boldsymbol{\beta}}{p-1}$ \\
Error & $\text{SSE}$ & $n-p$ & $\text{MSE} = \frac{\text{SSE}}{n-p}$ & $\sigma^2$ \\
Total & $\text{SST}$ & $n-1$ & $\text{MST} = \frac{\text{SST}}{n-1}$ & \\
\bottomrule
\end{tabular}
\end{table}

\textbf{Interpretación Geométrica}

En el espacio vectorial $\mathbb{R}^n$:

\begin{itemize}
\item $\mathbf{Y}$: Vector de observaciones
\item $\mathbf{\hat{Y}}$: Proyección de $\mathbf{Y}$ en el espacio columna de $\mathbf{X}$
\item $\boldsymbol{\varepsilon}$: Vector residual (ortogonal al espacio columna de $\mathbf{X}$)
\end{itemize}

La descomposición corresponde al teorema de Pitágoras:

\[
\|\mathbf{Y} - \bar{y}\mathbf{1}\|^2 = \|\mathbf{\hat{Y}} - \bar{y}\mathbf{1}\|^2 + \|\mathbf{Y} - \mathbf{\hat{Y}}\|^2
\]

\subsection{Propiedades Estadísticas}

\subsubsection{Distribuciones bajo Normalidad}

Si $\varepsilon_i \sim N(0, \sigma^2)$:

\begin{align*}
\frac{\text{SSE}}{\sigma^2} &\sim \chi^2(n-p) \\
\frac{\text{SSR}}{\sigma^2} &\sim \chi^2(p-1, \lambda) \quad \text{(no central)} \\
\frac{\text{SST}}{\sigma^2} &\sim \chi^2(n-1, \lambda')
\end{align*}

donde $\lambda = \frac{\boldsymbol{\beta}^T\mathbf{X}^T\mathbf{X}\boldsymbol{\beta}}{2\sigma^2}$ es el parámetro de no centralidad.

\subsubsection{Independencia Estadística}

\begin{itemize}
\item $\text{SSR}$ y $\text{SSE}$ son independientes
\item $\text{SSE}$ es independiente de $\hat{\boldsymbol{\beta}}$
\item $\text{SSR}$ y $\text{SST}$ no son independientes
\end{itemize}

\subsection{Aplicaciones en Inferencia}

\subsubsection{Prueba F de Significancia Global}

\[
F = \frac{\text{MSR}}{\text{MSE}} = \frac{\text{SSR}/(p-1)}{\text{SSE}/(n-p)} \sim F(p-1, n-p)
\]

Bajo $H_0: \beta_1 = \beta_2 = \dots = \beta_{p-1} = 0$

\subsubsection{Estimación de $\sigma^2$}

\[
\hat{\sigma}^2 = \text{MSE} = \frac{\text{SSE}}{n-p}
\]

es un estimador insesgado de $\sigma^2$.

\subsection{Ejemplo Numérico Detallado}

Para un modelo con $n=10$, $p=3$:

\begin{align*}
\text{SST} &= 150.0 \\
\text{SSR} &= 105.0 \\
\text{SSE} &= 45.0 \\
R^2 &= \frac{105.0}{150.0} = 0.70 \\
\text{MSE} &= \frac{45.0}{7} = 6.4286 \\
F &= \frac{105.0/2}{45.0/7} = \frac{52.5}{6.4286} = 8.1667
\end{align*}


%%%%%%%%%%%%%%%%%%%%%%%%%%%%%%%%%%%%%%%%%%%%%%%%%%%%%%%%%%%%%%%%
\section{Métricas de Bondad de Ajuste}

\subsection{Coeficiente de Determinación ($R^2$)}

\[
R^2 = 1 - \frac{\text{SSE}}{\text{SST}} = \frac{\text{SSR}}{\text{SST}}
\]

Interpretación: Proporción de la varianza en la variable dependiente que es explicada por el modelo.

\subsection{$R^2$ Ajustado}

\[
R^2_{\text{adj}} = 1 - \left(\frac{n-1}{n-p}\right)(1 - R^2)
\]

donde $p$ es el número de parámetros (incluyendo el intercepto).


\subsection{Error Cuadrático Medio (MSE)}

\[
\text{MSE} = \frac{1}{n}\sum_{i=1}^{n}(y_i - \hat{y}_i)^2 = \frac{\text{SSE}}{n}
\]

\subsection{Raíz del Error Cuadrático Medio (RMSE)}

\[
\text{RMSE} = \sqrt{\text{MSE}}
\]

\subsection{Error Absoluto Medio (MAE)}

\[
\text{MAE} = \frac{1}{n}\sum_{i=1}^{n}|y_i - \hat{y}_i|
\]

\section{Propiedades y Limitaciones}

\subsection{Propiedades del $R^2$}

\begin{enumerate}
\item $0 \leq R^2 \leq 1$
\item $R^2 = 1$ cuando SSE = 0 (ajuste perfecto)
\item $R^2 = 0$ cuando SSR = 0 (el modelo no explica nada)
\item El $R^2$ siempre aumenta al añadir más variables
\end{enumerate}

\subsection{Limitaciones del $R^2$}

\begin{itemize}
\item No indica si el modelo está correctamente especificado
\item No detecta relaciones no lineales
\item Puede ser inflado artificialmente añadiendo variables irrelevantes
\item No considera la significancia estadística de los coeficientes
\end{itemize}

\section{Evaluación Estadística}

\subsection{Prueba F de significancia global}

\[
F = \frac{\text{SSR}/(p-1)}{\text{SSE}/(n-p)} \sim F(p-1, n-p)
\]

Hipótesis:
\begin{itemize}
\item $H_0: \beta_1 = \beta_2 = \dots = \beta_{p-1} = 0$
\item $H_1: \text{Al menos un } \beta_j \neq 0$
\end{itemize}

\subsection{Error Estándar de los Residuales}

\[
s = \sqrt{\frac{\text{SSE}}{n-p}}
\]

\section{Ejemplo Numérico}

\begin{table}[h]
\centering
\caption{Descomposición de la variabilidad}
\begin{tabular}{lccc}
\toprule
Fuente & Suma de Cuadrados & Grados de Libertad & Cuadrado Medio \\
\midrule
Regresión & SSR & $p-1$ & MSR = SSR/(p-1) \\
Error & SSE & $n-p$ & MSE = SSE/(n-p) \\
Total & SST & $n-1$ & \\
\bottomrule
\end{tabular}
\end{table}

\section{Interpretación Práctica}

\subsection{Valores de Referencia}

\begin{itemize}
\item $R^2 < 0.3$: Ajuste pobre
\item $0.3 \leq R^2 < 0.7$: Ajuste moderado
\item $R^2 \geq 0.7$: Ajuste bueno
\end{itemize}

\subsection{Consideraciones Importantes}

\begin{enumerate}
\item Comparar $R^2$ entre modelos con la misma variable dependiente
\item Preferir $R^2$ ajustado cuando los modelos tienen diferente número de predictores
\item Considerar el contexto del problema y la variabilidad inherente de los datos
\item Validar con métricas complementarias (MSE, MAE)
\end{enumerate}

\section{Conclusiones}

La evaluación integral es necesaria porque diferentes métricas capturan aspectos distintos del desempeño del modelo:

\begin{align*}
R^2 &= 1 - \frac{\text{SSE}}{\text{SST}} \quad \text{(Varianza explicada)} \\
\text{MSE} &= \frac{1}{n}\sum_{i=1}^n (y_i - \hat{y}_i)^2 \quad \text{(Error cuadrático)} \\
\text{MAE} &= \frac{1}{n}\sum_{i=1}^n |y_i - \hat{y}_i| \quad \text{(Error absoluto)}
\end{align*}

Cada métrica proporciona información diferente sobre el comportamiento del modelo. La necesidad de múltiples métricas surge de sus diferentes propiedades matemáticas:

\begin{itemize}
\item \textbf{$R^2$:} Es adimensional pero susceptible a overfitting
\[
R^2_{\text{adj}} = 1 - \left(\frac{n-1}{n-p}\right)(1 - R^2)
\]
\item \textbf{MSE:} Penaliza más los errores grandes (diferenciable)
\[
\frac{\partial \text{MSE}}{\partial \hat{y}_i} = \frac{2}{n}(y_i - \hat{y}_i)
\]
\item \textbf{MAE:} Más robusto a outliers pero no diferenciable en cero
\[
\frac{\partial \text{MAE}}{\partial \hat{y}_i} = \text{sign}(y_i - \hat{y}_i)
\end{itemize}

\subsection{Considerar la significancia estadística}

La significancia estadística se verifica mediante:

\begin{align*}
F\text{-test:} &\quad F = \frac{\text{MSR}}{\text{MSE}} \sim F(p-1, n-p) \\
t\text{-test:} &\quad t_j = \frac{\hat{\beta}_j}{\text{SE}(\hat{\beta}_j)} \sim t(n-p)
\end{align*}

Un $R^2$ alto con p-values no significativos indica posible overfitting.

\subsection{Validar con datos de prueba}

La validación externa es crucial porque:

\[
\mathbb{E}[\text{MSE}_{\text{entrenamiento}}] \leq \mathbb{E}[\text{MSE}_{\text{prueba}}]
\]

El error de generalización se puede estimar mediante:

\[
\text{MSE}_{\text{prueba}} = \frac{1}{m}\sum_{i=1}^m (y_i^{\text{(test)}} - \hat{y}_i^{\text{(test)}})^2
\]

\subsection{Evaluar supuestos del modelo}

Los supuestos del modelo tienen implicaciones matemáticas precisas:

\subsubsection{Linealidad}
\[
\mathbb{E}[\varepsilon|\mathbf{X}] = 0 \quad \text{(El valor esperado de los errores es cero)}
\]

\subsubsection{Homocedasticidad}
\[
\text{Var}(\varepsilon|\mathbf{X}) = \sigma^2\mathbf{I} \quad \text{(Varianza constante)}
\]

\subsubsection{Independencia}
\[
\mathbb{E}[\varepsilon_i\varepsilon_j] = 0 \quad \text{para } i \neq j
\]

\subsubsection{Normalidad}
\[
\varepsilon \sim N(0, \sigma^2\mathbf{I}) \quad \text{(Para inferencia exacta)}
\]

\subsection{Un alto $R^2$ no garantiza un buen modelo predictivo}

Matemáticamente, un $R^2$ alto puede coexistir con:

\begin{align*}
\text{Overfitting:} &\quad R^2_{\text{entrenamiento}} \gg R^2_{\text{prueba}} \\
\text{Autocorrelación:} &\quad \mathbb{E}[\varepsilon_i\varepsilon_{i-1}] \neq 0 \\
\text{Heterocedasticidad:} &\quad \text{Var}(\varepsilon_i) \neq \sigma^2
\end{align*}

\subsection{Implicaciones para la predicción}

La calidad predictiva requiere:

\[
\mathbb{E}[(\hat{Y} - Y)^2] = \text{Var}(\hat{Y}) + \text{Bias}(\hat{Y})^2 + \text{Var}(\varepsilon)
\]

donde:
\begin{itemize}
\item $\text{Var}(\hat{Y})$: Varianza del estimador
\item $\text{Bias}(\hat{Y})$: Sesgo del estimador
\item $\text{Var}(\varepsilon)$: Varianza irreducible
\end{itemize}

\subsection{Recomendaciones Basadas en Evidencia Matemática}

\begin{enumerate}
\item \textbf{Validación Cruzada:}
\[
\text{CV}(k) = \frac{1}{k}\sum_{i=1}^k \text{MSE}_{\text{fold } i}
\]

\item \textbf{Regularización} (para mejorar generalización):
\[
\hat{\boldsymbol{\beta}}_{\text{ridge}} = \arg\min_{\boldsymbol{\beta}} \left\{\text{SSE} + \lambda\sum_{j=1}^p \beta_j^2\right\}
\]

\item \textbf{Análisis de Residuales:}
\[
r_i = y_i - \hat{y}_i \quad \text{deben ser } \sim N(0, \sigma^2)
\]

\item \textbf{Pruebas de Especificación:}
\begin{align*}
\text{Durbin-Watson:} &\quad d = \frac{\sum_{t=2}^n (r_t - r_{t-1})^2}{\sum_{t=1}^n r_t^2} \\
\text{Breusch-Pagan:} &\quad \text{Para heterocedasticidad}
\end{align*}
\end{enumerate}

\subsection{Conclusiones Finales}

La evaluación de modelos requiere un enfoque multifacético que combine:

\begin{itemize}
\item \textbf{Métricas cuantitativas} ($R^2$, MSE, MAE)
\item \textbf{Pruebas estadísticas} (F-test, t-test, pruebas de supuestos)
\item \textbf{Validación empírica} (datos de prueba, validación cruzada)
\item \textbf{Análisis cualitativo} (interpretabilidad, relevancia práctica)
\end{itemize}

Solo mediante esta aproximación integral se puede garantizar que un modelo con alta bondad de ajuste en muestra también tendrá buen desempeño predictivo en situaciones reales.


\end{document}
